%%%%%%%%%%%%%%%%%%%%%%%%%%%%%%%%%%%%%%%%%%%%%%%%%%%%%%%%%%%%
% 9.4 PARTIAL DIFFERENTIAL EQUATIONS
% The Composition of Advanced Mathematics
%%%%%%%%%%%%%%%%%%%%%%%%%%%%%%%%%%%%%%%%%%%%%%%%%%%%%%%%%%%%

\section{Partial Differential Equations}
\label{sec:partial-differential-equations}

\prerequisite{
Multivariable calculus, partial derivatives, multiple integrals, vector
calculus, ordinary differential equations, Fourier series, transform
methods, linear algebra, eigenvalues and eigenvectors, and basic boundary-
value problems.
}

\learningobjectives{
By the end of this section, the reader should be able to:
\begin{itemize}
    \item define a partial differential equation;
    \item determine the order of a PDE;
    \item distinguish linear, semilinear, quasilinear, and nonlinear PDEs;
    \item distinguish homogeneous and nonhomogeneous PDEs;
    \item identify boundary and initial conditions;
    \item distinguish Dirichlet, Neumann, and Robin boundary conditions;
    \item classify second-order linear PDEs as elliptic, parabolic, or
          hyperbolic;
    \item recognize Laplace's equation;
    \item recognize Poisson's equation;
    \item recognize the heat equation;
    \item recognize the wave equation;
    \item understand the transport equation;
    \item solve elementary first-order transport equations;
    \item understand the method of characteristics;
    \item use separation of variables;
    \item formulate eigenvalue problems arising from separation;
    \item understand Fourier-series solutions of boundary-value problems;
    \item solve the one-dimensional heat equation on a finite interval;
    \item solve the one-dimensional wave equation on a finite interval;
    \item understand steady-state temperature problems;
    \item solve basic Laplace-equation problems using separation of
          variables;
    \item understand harmonic functions;
    \item recognize the maximum principle;
    \item understand uniqueness conceptually;
    \item understand energy methods conceptually;
    \item understand d'Alembert's formula for the wave equation;
    \item understand diffusion and smoothing;
    \item distinguish finite and infinite propagation speed;
    \item understand fundamental solutions conceptually;
    \item recognize Green's functions;
    \item understand transform methods for PDEs;
    \item recognize nonlinear PDEs;
    \item understand conservation laws conceptually;
    \item recognize shocks and weak solutions conceptually;
    \item understand numerical PDE methods conceptually;
    \item connect PDEs with physics, engineering, geometry, probability,
          and numerical mathematics.
\end{itemize}
}

%%%%%%%%%%%%%%%%%%%%%%%%%%%%%%%%%%%%%%%%%%%%%%%%%%%%%%%%%%%%
\subsection{What Is a Partial Differential Equation?}
\label{subsec:what-is-pde}
%%%%%%%%%%%%%%%%%%%%%%%%%%%%%%%%%%%%%%%%%%%%%%%%%%%%%%%%%%%%

A partial differential equation, abbreviated PDE, is an equation involving
an unknown function of several independent variables and its partial
derivatives.

For example,

\[
\boxed{
\frac{\partial u}{\partial t}
=
k
\frac{\partial^2u}{\partial x^2}
}
\]

is a PDE for the function

\[
u=u(x,t).
\]

Another example is

\[
\boxed{
u_{xx}+u_{yy}=0.
}
\]

%%%%%%%%%%%%%%%%%%%%%%%%%%%%%%%%%%%%%%%%%%%%%%%%%%%%%%%%%%%%
\subsection{Ordinary Versus Partial Differential Equations}
\label{subsec:ode-versus-pde-extended}
%%%%%%%%%%%%%%%%%%%%%%%%%%%%%%%%%%%%%%%%%%%%%%%%%%%%%%%%%%%%

An ordinary differential equation involves one independent variable.

For example,

\[
\boxed{
y''+y=0.
}
\]

A partial differential equation involves two or more independent
variables.

For example,

\[
\boxed{
u_t
=
ku_{xx}.
}
\]

The unknown function may depend on variables such as:

\[
\boxed{
x,\ y,\ z,\ t.
}
\]

%%%%%%%%%%%%%%%%%%%%%%%%%%%%%%%%%%%%%%%%%%%%%%%%%%%%%%%%%%%%
\subsection{General Form of a PDE}
\label{subsec:general-form-pde}
%%%%%%%%%%%%%%%%%%%%%%%%%%%%%%%%%%%%%%%%%%%%%%%%%%%%%%%%%%%%

A PDE may be written abstractly as

\[
\boxed{
F
\left(
x_1,\ldots,x_n,
u,
Du,
D^2u,
\ldots,
D^mu
\right)
=
0.
}
\]

Here:

\[
Du
\]

represents first partial derivatives,

\[
D^2u
\]

represents second partial derivatives, and so on.

%%%%%%%%%%%%%%%%%%%%%%%%%%%%%%%%%%%%%%%%%%%%%%%%%%%%%%%%%%%%
\subsection{Order of a PDE}
\label{subsec:pde-order}
%%%%%%%%%%%%%%%%%%%%%%%%%%%%%%%%%%%%%%%%%%%%%%%%%%%%%%%%%%%%

\begin{definitionbox}{Order of a PDE}
The order of a partial differential equation is the highest order of any
partial derivative appearing in the equation.
\end{definitionbox}

For example,

\[
u_{xxx}
+
u_{xy}
-
u
=
0
\]

is third order.

Thus,

\[
\boxed{
\operatorname{order}=3.
}
\]

%%%%%%%%%%%%%%%%%%%%%%%%%%%%%%%%%%%%%%%%%%%%%%%%%%%%%%%%%%%%
\subsection{Linear PDEs}
\label{subsec:linear-pdes}
%%%%%%%%%%%%%%%%%%%%%%%%%%%%%%%%%%%%%%%%%%%%%%%%%%%%%%%%%%%%

A PDE is linear if the unknown function and all its derivatives appear
linearly.

A second-order linear PDE in two variables may be written

\[
\boxed{
A u_{xx}
+
2B u_{xy}
+
C u_{yy}
+
D u_x
+
E u_y
+
F u
=
G,
}
\]

where the coefficients may depend on the independent variables but not on
\(u\) or its derivatives.

%%%%%%%%%%%%%%%%%%%%%%%%%%%%%%%%%%%%%%%%%%%%%%%%%%%%%%%%%%%%
\subsection{Homogeneous Linear PDEs}
\label{subsec:homogeneous-linear-pdes}
%%%%%%%%%%%%%%%%%%%%%%%%%%%%%%%%%%%%%%%%%%%%%%%%%%%%%%%%%%%%

A linear PDE is homogeneous when the forcing term is zero.

For example,

\[
\boxed{
u_t-ku_{xx}=0.
}
\]

Similarly,

\[
\boxed{
u_{xx}+u_{yy}=0
}
\]

is homogeneous.

%%%%%%%%%%%%%%%%%%%%%%%%%%%%%%%%%%%%%%%%%%%%%%%%%%%%%%%%%%%%
\subsection{Nonhomogeneous Linear PDEs}
\label{subsec:nonhomogeneous-linear-pdes}
%%%%%%%%%%%%%%%%%%%%%%%%%%%%%%%%%%%%%%%%%%%%%%%%%%%%%%%%%%%%

A linear PDE is nonhomogeneous when

\[
\boxed{
G\neq0.
}
\]

For example,

\[
\boxed{
u_{xx}+u_{yy}
=
f(x,y).
}
\]

The right side acts as a source or forcing term.

%%%%%%%%%%%%%%%%%%%%%%%%%%%%%%%%%%%%%%%%%%%%%%%%%%%%%%%%%%%%
\subsection{Semilinear PDEs}
\label{subsec:semilinear-pdes}
%%%%%%%%%%%%%%%%%%%%%%%%%%%%%%%%%%%%%%%%%%%%%%%%%%%%%%%%%%%%

A PDE is semilinear when the highest-order derivatives appear linearly,
while lower-order terms may depend nonlinearly on \(u\).

For example,

\[
\boxed{
u_t-u_{xx}
=
u-u^3.
}
\]

The second-order derivative remains linear.

%%%%%%%%%%%%%%%%%%%%%%%%%%%%%%%%%%%%%%%%%%%%%%%%%%%%%%%%%%%%
\subsection{Quasilinear PDEs}
\label{subsec:quasilinear-pdes}
%%%%%%%%%%%%%%%%%%%%%%%%%%%%%%%%%%%%%%%%%%%%%%%%%%%%%%%%%%%%

A PDE is quasilinear when the highest-order derivatives appear linearly,
but their coefficients may depend on the unknown function or lower-order
derivatives.

For example,

\[
\boxed{
u_t
+
u\,u_x
=
0
}
\]

is first-order quasilinear.

%%%%%%%%%%%%%%%%%%%%%%%%%%%%%%%%%%%%%%%%%%%%%%%%%%%%%%%%%%%%
\subsection{Fully Nonlinear PDEs}
\label{subsec:fully-nonlinear-pdes}
%%%%%%%%%%%%%%%%%%%%%%%%%%%%%%%%%%%%%%%%%%%%%%%%%%%%%%%%%%%%

A PDE is fully nonlinear when the highest-order derivatives themselves
appear nonlinearly.

For example,

\[
\boxed{
u_{xx}u_{yy}
-
u_{xy}^2
=
1.
}
\]

This is a form of the Monge--Ampère equation.

%%%%%%%%%%%%%%%%%%%%%%%%%%%%%%%%%%%%%%%%%%%%%%%%%%%%%%%%%%%%
\subsection{Initial Conditions}
\label{subsec:pde-initial-conditions}
%%%%%%%%%%%%%%%%%%%%%%%%%%%%%%%%%%%%%%%%%%%%%%%%%%%%%%%%%%%%

A time-dependent PDE often requires initial data.

For example,

\[
\boxed{
u(x,0)
=
f(x).
}
\]

For a second-order time equation such as the wave equation, one may need

\[
\boxed{
u(x,0)=f(x),
}
\]

and

\[
\boxed{
u_t(x,0)=g(x).
}
\]

%%%%%%%%%%%%%%%%%%%%%%%%%%%%%%%%%%%%%%%%%%%%%%%%%%%%%%%%%%%%
\subsection{Boundary Conditions}
\label{subsec:pde-boundary-conditions}
%%%%%%%%%%%%%%%%%%%%%%%%%%%%%%%%%%%%%%%%%%%%%%%%%%%%%%%%%%%%

PDEs posed on a spatial domain usually require conditions along the
boundary.

Common types include:

\[
\boxed{
\text{Dirichlet},
}
\]

\[
\boxed{
\text{Neumann},
}
\]

and

\[
\boxed{
\text{Robin}.
}
\]

%%%%%%%%%%%%%%%%%%%%%%%%%%%%%%%%%%%%%%%%%%%%%%%%%%%%%%%%%%%%
\subsection{Dirichlet Boundary Conditions}
\label{subsec:dirichlet-boundary-conditions}
%%%%%%%%%%%%%%%%%%%%%%%%%%%%%%%%%%%%%%%%%%%%%%%%%%%%%%%%%%%%

Dirichlet conditions specify the value of the unknown function on the
boundary.

For example,

\[
\boxed{
u(0,t)=0,
\qquad
u(L,t)=0.
}
\]

In a heat problem, these may represent fixed temperatures at the ends of
a rod.

%%%%%%%%%%%%%%%%%%%%%%%%%%%%%%%%%%%%%%%%%%%%%%%%%%%%%%%%%%%%
\subsection{Neumann Boundary Conditions}
\label{subsec:neumann-boundary-conditions}
%%%%%%%%%%%%%%%%%%%%%%%%%%%%%%%%%%%%%%%%%%%%%%%%%%%%%%%%%%%%

Neumann conditions specify a normal derivative on the boundary.

For a one-dimensional interval,

\[
\boxed{
u_x(0,t)=0,
\qquad
u_x(L,t)=0.
}
\]

In heat conduction, zero derivative can represent an insulated boundary.

%%%%%%%%%%%%%%%%%%%%%%%%%%%%%%%%%%%%%%%%%%%%%%%%%%%%%%%%%%%%
\subsection{Robin Boundary Conditions}
\label{subsec:robin-boundary-conditions}
%%%%%%%%%%%%%%%%%%%%%%%%%%%%%%%%%%%%%%%%%%%%%%%%%%%%%%%%%%%%

Robin conditions combine function values and derivatives:

\[
\boxed{
a u
+
b\frac{\partial u}{\partial n}
=
c.
}
\]

They arise naturally in convection, heat-transfer, and impedance-type
boundary models.

%%%%%%%%%%%%%%%%%%%%%%%%%%%%%%%%%%%%%%%%%%%%%%%%%%%%%%%%%%%%
\subsection{Mixed Boundary Conditions}
\label{subsec:mixed-boundary-conditions}
%%%%%%%%%%%%%%%%%%%%%%%%%%%%%%%%%%%%%%%%%%%%%%%%%%%%%%%%%%%%

Different parts of the boundary may satisfy different types of
conditions.

For example,

\[
u(0,t)=0,
\]

while

\[
u_x(L,t)=0.
\]

Such problems are called mixed boundary-value problems.

%%%%%%%%%%%%%%%%%%%%%%%%%%%%%%%%%%%%%%%%%%%%%%%%%%%%%%%%%%%%
\subsection{Classification of Second-Order PDEs}
\label{subsec:second-order-pde-classification}
%%%%%%%%%%%%%%%%%%%%%%%%%%%%%%%%%%%%%%%%%%%%%%%%%%%%%%%%%%%%

Consider the principal second-order part

\[
\boxed{
A u_{xx}
+
2B u_{xy}
+
C u_{yy}.
}
\]

The discriminant is

\[
\boxed{
B^2-AC.
}
\]

At a point, the PDE is classified as:

\[
\boxed{
B^2-AC<0
\Rightarrow
\text{elliptic},
}
\]

\[
\boxed{
B^2-AC=0
\Rightarrow
\text{parabolic},
}
\]

and

\[
\boxed{
B^2-AC>0
\Rightarrow
\text{hyperbolic}.
}
\]

%%%%%%%%%%%%%%%%%%%%%%%%%%%%%%%%%%%%%%%%%%%%%%%%%%%%%%%%%%%%
\subsection{Elliptic Equations}
\label{subsec:elliptic-equations}
%%%%%%%%%%%%%%%%%%%%%%%%%%%%%%%%%%%%%%%%%%%%%%%%%%%%%%%%%%%%

Elliptic equations often describe steady-state or equilibrium phenomena.

The fundamental example is Laplace's equation:

\[
\boxed{
\Delta u=0.
}
\]

In two dimensions,

\[
\boxed{
u_{xx}+u_{yy}=0.
}
\]

%%%%%%%%%%%%%%%%%%%%%%%%%%%%%%%%%%%%%%%%%%%%%%%%%%%%%%%%%%%%
\subsection{Parabolic Equations}
\label{subsec:parabolic-equations}
%%%%%%%%%%%%%%%%%%%%%%%%%%%%%%%%%%%%%%%%%%%%%%%%%%%%%%%%%%%%

Parabolic equations commonly model diffusion and smoothing.

The standard example is the heat equation:

\[
\boxed{
u_t
=
\alpha^2u_{xx}.
}
\]

Here,

\[
\alpha^2>0
\]

is a diffusion coefficient.

%%%%%%%%%%%%%%%%%%%%%%%%%%%%%%%%%%%%%%%%%%%%%%%%%%%%%%%%%%%%
\subsection{Hyperbolic Equations}
\label{subsec:hyperbolic-equations}
%%%%%%%%%%%%%%%%%%%%%%%%%%%%%%%%%%%%%%%%%%%%%%%%%%%%%%%%%%%%

Hyperbolic equations often describe wave propagation and transport.

The standard wave equation is

\[
\boxed{
u_{tt}
=
c^2u_{xx}.
}
\]

The parameter

\[
c>0
\]

is the wave speed.

%%%%%%%%%%%%%%%%%%%%%%%%%%%%%%%%%%%%%%%%%%%%%%%%%%%%%%%%%%%%
\subsection{The Laplacian}
\label{subsec:laplacian}
%%%%%%%%%%%%%%%%%%%%%%%%%%%%%%%%%%%%%%%%%%%%%%%%%%%%%%%%%%%%

For a function

\[
u=u(x_1,\ldots,x_n),
\]

the Laplacian is

\[
\boxed{
\Delta u
=
\sum_{j=1}^{n}
\frac{\partial^2u}{\partial x_j^2}.
}
\]

In two dimensions,

\[
\boxed{
\Delta u
=
u_{xx}+u_{yy}.
}
\]

In three dimensions,

\[
\boxed{
\Delta u
=
u_{xx}+u_{yy}+u_{zz}.
}
\]

%%%%%%%%%%%%%%%%%%%%%%%%%%%%%%%%%%%%%%%%%%%%%%%%%%%%%%%%%%%%
\subsection{Laplace's Equation}
\label{subsec:laplaces-equation}
%%%%%%%%%%%%%%%%%%%%%%%%%%%%%%%%%%%%%%%%%%%%%%%%%%%%%%%%%%%%

Laplace's equation is

\[
\boxed{
\Delta u=0.
}
\]

Solutions are called harmonic functions.

Applications include:

\[
\boxed{
\text{steady-state temperature},
}
\]

\[
\boxed{
\text{electrostatic potential},
}
\]

\[
\boxed{
\text{gravitational potential},
}
\]

and

\[
\boxed{
\text{incompressible potential flow}.
}
\]

%%%%%%%%%%%%%%%%%%%%%%%%%%%%%%%%%%%%%%%%%%%%%%%%%%%%%%%%%%%%
\subsection{Poisson's Equation}
\label{subsec:poissons-equation}
%%%%%%%%%%%%%%%%%%%%%%%%%%%%%%%%%%%%%%%%%%%%%%%%%%%%%%%%%%%%

Poisson's equation is

\[
\boxed{
\Delta u
=
f.
}
\]

The function

\[
f
\]

represents a source distribution.

Laplace's equation is the special case

\[
\boxed{
f=0.
}
\]

%%%%%%%%%%%%%%%%%%%%%%%%%%%%%%%%%%%%%%%%%%%%%%%%%%%%%%%%%%%%
\subsection{Heat Equation}
\label{subsec:heat-equation}
%%%%%%%%%%%%%%%%%%%%%%%%%%%%%%%%%%%%%%%%%%%%%%%%%%%%%%%%%%%%

The one-dimensional heat equation is

\[
\boxed{
u_t
=
\alpha^2u_{xx}.
}
\]

Here

\[
u(x,t)
\]

may represent temperature at position \(x\) and time \(t\).

The equation says temporal change is proportional to spatial curvature.

%%%%%%%%%%%%%%%%%%%%%%%%%%%%%%%%%%%%%%%%%%%%%%%%%%%%%%%%%%%%
\subsection{Physical Interpretation of Heat Diffusion}
\label{subsec:heat-physical-interpretation}
%%%%%%%%%%%%%%%%%%%%%%%%%%%%%%%%%%%%%%%%%%%%%%%%%%%%%%%%%%%%

If

\[
u_{xx}>0,
\]

the temperature profile is locally concave upward, so nearby temperatures
tend to raise the local value.

If

\[
u_{xx}<0,
\]

the local point tends to cool.

Thus diffusion smooths spatial variation.

%%%%%%%%%%%%%%%%%%%%%%%%%%%%%%%%%%%%%%%%%%%%%%%%%%%%%%%%%%%%
\subsection{Wave Equation}
\label{subsec:wave-equation}
%%%%%%%%%%%%%%%%%%%%%%%%%%%%%%%%%%%%%%%%%%%%%%%%%%%%%%%%%%%%

The one-dimensional wave equation is

\[
\boxed{
u_{tt}
=
c^2u_{xx}.
}
\]

It models vibrating strings, acoustic waves, elastic waves, and other
propagating disturbances.

Unlike the heat equation, it contains a second derivative in time.

%%%%%%%%%%%%%%%%%%%%%%%%%%%%%%%%%%%%%%%%%%%%%%%%%%%%%%%%%%%%
\subsection{Transport Equation}
\label{subsec:transport-equation}
%%%%%%%%%%%%%%%%%%%%%%%%%%%%%%%%%%%%%%%%%%%%%%%%%%%%%%%%%%%%

A basic transport equation is

\[
\boxed{
u_t
+
c u_x
=
0.
}
\]

It represents a profile moving with constant velocity \(c\).

If

\[
u(x,0)
=
f(x),
\]

then the solution is

\[
\boxed{
u(x,t)
=
f(x-ct).
}
\]

%%%%%%%%%%%%%%%%%%%%%%%%%%%%%%%%%%%%%%%%%%%%%%%%%%%%%%%%%%%%
\subsection{Worked Example: Transport Equation}
\label{subsec:worked-transport-equation}
%%%%%%%%%%%%%%%%%%%%%%%%%%%%%%%%%%%%%%%%%%%%%%%%%%%%%%%%%%%%

\begin{workedexamplebox}{Moving Profile}

Solve

\[
u_t+3u_x=0,
\qquad
u(x,0)=e^{-x^2}.
\]

The transport speed is

\[
c=3.
\]

Therefore the initial profile moves to the right without changing shape:

\begin{answerbox}
\[
\boxed{
u(x,t)
=
e^{-(x-3t)^2}.
}
\]
\end{answerbox}

\end{workedexamplebox}

%%%%%%%%%%%%%%%%%%%%%%%%%%%%%%%%%%%%%%%%%%%%%%%%%%%%%%%%%%%%
\subsection{Method of Characteristics}
\label{subsec:method-of-characteristics}
%%%%%%%%%%%%%%%%%%%%%%%%%%%%%%%%%%%%%%%%%%%%%%%%%%%%%%%%%%%%

The method of characteristics converts certain first-order PDEs into ODEs
along special curves.

Consider

\[
\boxed{
a(x,t,u)u_x
+
b(x,t,u)u_t
=
c(x,t,u).
}
\]

Characteristic curves satisfy a system such as

\[
\boxed{
\frac{dx}{ds}=a,
\qquad
\frac{dt}{ds}=b,
\qquad
\frac{du}{ds}=c.
}
\]

Along these curves, the PDE becomes an ordinary differential equation.

%%%%%%%%%%%%%%%%%%%%%%%%%%%%%%%%%%%%%%%%%%%%%%%%%%%%%%%%%%%%
\subsection{Characteristics for Constant Transport}
\label{subsec:transport-characteristics}
%%%%%%%%%%%%%%%%%%%%%%%%%%%%%%%%%%%%%%%%%%%%%%%%%%%%%%%%%%%%

For

\[
u_t+cu_x=0,
\]

the characteristic equations are

\[
\boxed{
\frac{dx}{dt}=c.
}
\]

Integrating,

\[
\boxed{
x-ct=C.
}
\]

Along each characteristic,

\[
\boxed{
\frac{du}{dt}=0.
}
\]

Thus \(u\) is constant on lines

\[
x-ct=C.
\]

%%%%%%%%%%%%%%%%%%%%%%%%%%%%%%%%%%%%%%%%%%%%%%%%%%%%%%%%%%%%
\subsection{Separation of Variables}
\label{subsec:pde-separation-variables}
%%%%%%%%%%%%%%%%%%%%%%%%%%%%%%%%%%%%%%%%%%%%%%%%%%%%%%%%%%%%

A major solution method assumes a product form

\[
\boxed{
u(x,t)
=
X(x)T(t).
}
\]

Substitution into a PDE separates the variables into ODEs.

This works especially well for:

\[
\boxed{
\text{linear homogeneous PDEs}
}
\]

with simple boundary conditions.

%%%%%%%%%%%%%%%%%%%%%%%%%%%%%%%%%%%%%%%%%%%%%%%%%%%%%%%%%%%%
\subsection{Separation Constant}
\label{subsec:pde-separation-constant}
%%%%%%%%%%%%%%%%%%%%%%%%%%%%%%%%%%%%%%%%%%%%%%%%%%%%%%%%%%%%

Suppose substitution produces

\[
\frac{T'}{\alpha^2T}
=
\frac{X''}{X}.
\]

The left side depends only on \(t\).

The right side depends only on \(x\).

Therefore both must equal a constant:

\[
\boxed{
-\lambda.
}
\]

This yields two ODEs.

%%%%%%%%%%%%%%%%%%%%%%%%%%%%%%%%%%%%%%%%%%%%%%%%%%%%%%%%%%%%
\subsection{Spatial Eigenvalue Problem}
\label{subsec:pde-spatial-eigenvalue-problem}
%%%%%%%%%%%%%%%%%%%%%%%%%%%%%%%%%%%%%%%%%%%%%%%%%%%%%%%%%%%%

For the heat equation with fixed-end conditions,

\[
u_t
=
\alpha^2u_{xx},
\]

\[
u(0,t)=0,
\qquad
u(L,t)=0,
\]

separation yields

\[
\boxed{
X''+\lambda X=0,
}
\]

with

\[
\boxed{
X(0)=0,
\qquad
X(L)=0.
}
\]

This is an eigenvalue problem.

%%%%%%%%%%%%%%%%%%%%%%%%%%%%%%%%%%%%%%%%%%%%%%%%%%%%%%%%%%%%
\subsection{Eigenvalues for Fixed Ends}
\label{subsec:fixed-end-eigenvalues}
%%%%%%%%%%%%%%%%%%%%%%%%%%%%%%%%%%%%%%%%%%%%%%%%%%%%%%%%%%%%

Nontrivial solutions exist when

\[
\boxed{
\lambda_n
=
\left(
\frac{n\pi}{L}
\right)^2,
\qquad
n=1,2,\ldots
}
\]

with eigenfunctions

\[
\boxed{
X_n(x)
=
\sin
\left(
\frac{n\pi x}{L}
\right).
}
\]

%%%%%%%%%%%%%%%%%%%%%%%%%%%%%%%%%%%%%%%%%%%%%%%%%%%%%%%%%%%%
\subsection{Orthogonality of Sine Eigenfunctions}
\label{subsec:sine-eigenfunction-orthogonality}
%%%%%%%%%%%%%%%%%%%%%%%%%%%%%%%%%%%%%%%%%%%%%%%%%%%%%%%%%%%%

For positive integers \(m\neq n\),

\[
\boxed{
\int_0^L
\sin
\left(
\frac{m\pi x}{L}
\right)
\sin
\left(
\frac{n\pi x}{L}
\right)
\,dx
=
0.
}
\]

For \(m=n\),

\[
\boxed{
\int_0^L
\sin^2
\left(
\frac{n\pi x}{L}
\right)
\,dx
=
\frac{L}{2}.
}
\]

This orthogonality determines Fourier coefficients.

%%%%%%%%%%%%%%%%%%%%%%%%%%%%%%%%%%%%%%%%%%%%%%%%%%%%%%%%%%%%
\subsection{One-Dimensional Heat Equation}
\label{subsec:one-dimensional-heat-equation}
%%%%%%%%%%%%%%%%%%%%%%%%%%%%%%%%%%%%%%%%%%%%%%%%%%%%%%%%%%%%

Consider

\[
\boxed{
u_t
=
\alpha^2u_{xx},
\qquad
0<x<L,
\quad
t>0,
}
\]

with

\[
\boxed{
u(0,t)=0,
\qquad
u(L,t)=0,
}
\]

and

\[
\boxed{
u(x,0)=f(x).
}
\]

%%%%%%%%%%%%%%%%%%%%%%%%%%%%%%%%%%%%%%%%%%%%%%%%%%%%%%%%%%%%
\subsection{Separated Heat Modes}
\label{subsec:heat-separated-modes}
%%%%%%%%%%%%%%%%%%%%%%%%%%%%%%%%%%%%%%%%%%%%%%%%%%%%%%%%%%%%

For eigenvalue

\[
\lambda_n
=
\left(
\frac{n\pi}{L}
\right)^2,
\]

the time equation is

\[
\boxed{
T_n'
+
\alpha^2\lambda_nT_n
=
0.
}
\]

Hence

\[
\boxed{
T_n(t)
=
e^{-\alpha^2\lambda_nt}.
}
\]

Each mode decays exponentially.

%%%%%%%%%%%%%%%%%%%%%%%%%%%%%%%%%%%%%%%%%%%%%%%%%%%%%%%%%%%%
\subsection{Heat Equation Fourier-Series Solution}
\label{subsec:heat-fourier-series-solution}
%%%%%%%%%%%%%%%%%%%%%%%%%%%%%%%%%%%%%%%%%%%%%%%%%%%%%%%%%%%%

The solution is

\[
\boxed{
u(x,t)
=
\sum_{n=1}^{\infty}
b_n
e^{
-\alpha^2
(n\pi/L)^2
t
}
\sin
\left(
\frac{n\pi x}{L}
\right).
}
\]

The coefficients are determined from the initial condition:

\[
\boxed{
b_n
=
\frac{2}{L}
\int_0^L
f(x)
\sin
\left(
\frac{n\pi x}{L}
\right)
\,dx.
}
\]

%%%%%%%%%%%%%%%%%%%%%%%%%%%%%%%%%%%%%%%%%%%%%%%%%%%%%%%%%%%%
\subsection{Worked Example: Single Heat Mode}
\label{subsec:worked-single-heat-mode}
%%%%%%%%%%%%%%%%%%%%%%%%%%%%%%%%%%%%%%%%%%%%%%%%%%%%%%%%%%%%

\begin{workedexamplebox}{Heat Equation with One Fourier Mode}

Suppose

\[
u_t=u_{xx},
\qquad
0<x<\pi,
\]

with

\[
u(0,t)=u(\pi,t)=0,
\]

and

\[
u(x,0)=3\sin(2x).
\]

The initial condition is already an eigenfunction with

\[
n=2.
\]

Thus,

\begin{answerbox}
\[
\boxed{
u(x,t)
=
3e^{-4t}\sin(2x).
}
\]
\end{answerbox}

\end{workedexamplebox}

%%%%%%%%%%%%%%%%%%%%%%%%%%%%%%%%%%%%%%%%%%%%%%%%%%%%%%%%%%%%
\subsection{Long-Time Behavior of the Heat Equation}
\label{subsec:heat-long-time}
%%%%%%%%%%%%%%%%%%%%%%%%%%%%%%%%%%%%%%%%%%%%%%%%%%%%%%%%%%%%

Each term contains

\[
e^{-\alpha^2(n\pi/L)^2t}.
\]

Therefore,

\[
\boxed{
u(x,t)
\to0
}
\]

as

\[
t\to\infty
\]

for the zero-boundary problem under suitable initial data.

Higher-frequency modes decay faster because their eigenvalues are larger.

%%%%%%%%%%%%%%%%%%%%%%%%%%%%%%%%%%%%%%%%%%%%%%%%%%%%%%%%%%%%
\subsection{Diffusive Smoothing}
\label{subsec:diffusive-smoothing}
%%%%%%%%%%%%%%%%%%%%%%%%%%%%%%%%%%%%%%%%%%%%%%%%%%%%%%%%%%%%

The heat equation strongly damps high-frequency spatial oscillations.

A mode with frequency \(n\) decays like

\[
\boxed{
e^{-Cn^2t}.
}
\]

Thus rough initial data become smoother for positive time.

This smoothing behavior is characteristic of parabolic equations.

%%%%%%%%%%%%%%%%%%%%%%%%%%%%%%%%%%%%%%%%%%%%%%%%%%%%%%%%%%%%
\subsection{Nonzero Boundary Conditions}
\label{subsec:nonzero-heat-boundary}
%%%%%%%%%%%%%%%%%%%%%%%%%%%%%%%%%%%%%%%%%%%%%%%%%%%%%%%%%%%%

Suppose

\[
u(0,t)=A,
\qquad
u(L,t)=B.
\]

A common method writes

\[
\boxed{
u(x,t)
=
v(x,t)+w(x),
}
\]

where \(w(x)\) is a steady-state solution satisfying the boundary
conditions.

Then \(v\) satisfies homogeneous boundary conditions.

%%%%%%%%%%%%%%%%%%%%%%%%%%%%%%%%%%%%%%%%%%%%%%%%%%%%%%%%%%%%
\subsection{Steady-State Heat Equation}
\label{subsec:steady-state-heat}
%%%%%%%%%%%%%%%%%%%%%%%%%%%%%%%%%%%%%%%%%%%%%%%%%%%%%%%%%%%%

At steady state,

\[
u_t=0.
\]

The heat equation becomes

\[
\boxed{
u_{xx}=0.
}
\]

Thus,

\[
\boxed{
u(x)
=
Ax+B.
}
\]

For endpoint temperatures

\[
u(0)=T_0,
\qquad
u(L)=T_L,
\]

the steady profile is

\[
\boxed{
u(x)
=
T_0
+
\frac{
T_L-T_0
}{
L
}x.
}
\]

%%%%%%%%%%%%%%%%%%%%%%%%%%%%%%%%%%%%%%%%%%%%%%%%%%%%%%%%%%%%
\subsection{One-Dimensional Wave Equation}
\label{subsec:one-dimensional-wave-equation}
%%%%%%%%%%%%%%%%%%%%%%%%%%%%%%%%%%%%%%%%%%%%%%%%%%%%%%%%%%%%

Consider

\[
\boxed{
u_{tt}
=
c^2u_{xx},
\qquad
0<x<L,
}
\]

with fixed ends:

\[
\boxed{
u(0,t)=0,
\qquad
u(L,t)=0.
}
\]

Initial conditions are

\[
\boxed{
u(x,0)=f(x),
}
\]

and

\[
\boxed{
u_t(x,0)=g(x).
}
\]

%%%%%%%%%%%%%%%%%%%%%%%%%%%%%%%%%%%%%%%%%%%%%%%%%%%%%%%%%%%%
\subsection{Separated Wave Modes}
\label{subsec:wave-separated-modes}
%%%%%%%%%%%%%%%%%%%%%%%%%%%%%%%%%%%%%%%%%%%%%%%%%%%%%%%%%%%%

Separation gives the same spatial eigenfunctions:

\[
\boxed{
X_n(x)
=
\sin
\left(
\frac{n\pi x}{L}
\right).
}
\]

The time equation becomes

\[
\boxed{
T_n''
+
c^2
\left(
\frac{n\pi}{L}
\right)^2
T_n
=
0.
}
\]

Hence,

\[
\boxed{
T_n(t)
=
A_n
\cos
\left(
\frac{cn\pi t}{L}
\right)
+
B_n
\sin
\left(
\frac{cn\pi t}{L}
\right).
}
\]

%%%%%%%%%%%%%%%%%%%%%%%%%%%%%%%%%%%%%%%%%%%%%%%%%%%%%%%%%%%%
\subsection{Wave Equation Series Solution}
\label{subsec:wave-series-solution}
%%%%%%%%%%%%%%%%%%%%%%%%%%%%%%%%%%%%%%%%%%%%%%%%%%%%%%%%%%%%

The fixed-end solution is

\[
\boxed{
u(x,t)
=
\sum_{n=1}^{\infty}
\left[
A_n
\cos
\left(
\frac{cn\pi t}{L}
\right)
+
B_n
\sin
\left(
\frac{cn\pi t}{L}
\right)
\right]
\sin
\left(
\frac{n\pi x}{L}
\right).
}
\]

%%%%%%%%%%%%%%%%%%%%%%%%%%%%%%%%%%%%%%%%%%%%%%%%%%%%%%%%%%%%
\subsection{Wave Initial Displacement Coefficients}
\label{subsec:wave-initial-displacement}
%%%%%%%%%%%%%%%%%%%%%%%%%%%%%%%%%%%%%%%%%%%%%%%%%%%%%%%%%%%%

From

\[
u(x,0)=f(x),
\]

we obtain

\[
\boxed{
f(x)
=
\sum_{n=1}^{\infty}
A_n
\sin
\left(
\frac{n\pi x}{L}
\right).
}
\]

Therefore,

\[
\boxed{
A_n
=
\frac{2}{L}
\int_0^L
f(x)
\sin
\left(
\frac{n\pi x}{L}
\right)
\,dx.
}
\]

%%%%%%%%%%%%%%%%%%%%%%%%%%%%%%%%%%%%%%%%%%%%%%%%%%%%%%%%%%%%
\subsection{Wave Initial Velocity Coefficients}
\label{subsec:wave-initial-velocity}
%%%%%%%%%%%%%%%%%%%%%%%%%%%%%%%%%%%%%%%%%%%%%%%%%%%%%%%%%%%%

Differentiate:

\[
u_t(x,0)
=
\sum_{n=1}^{\infty}
\frac{cn\pi}{L}
B_n
\sin
\left(
\frac{n\pi x}{L}
\right).
\]

Since this equals \(g(x)\),

\[
\boxed{
B_n
=
\frac{2}{cn\pi}
\int_0^L
g(x)
\sin
\left(
\frac{n\pi x}{L}
\right)
\,dx.
}
\]

%%%%%%%%%%%%%%%%%%%%%%%%%%%%%%%%%%%%%%%%%%%%%%%%%%%%%%%%%%%%
\subsection{Worked Example: Single Wave Mode}
\label{subsec:worked-single-wave-mode}
%%%%%%%%%%%%%%%%%%%%%%%%%%%%%%%%%%%%%%%%%%%%%%%%%%%%%%%%%%%%

\begin{workedexamplebox}{Vibrating String Mode}

Suppose

\[
u_{tt}
=
4u_{xx},
\qquad
0<x<\pi,
\]

with

\[
u(0,t)=u(\pi,t)=0,
\]

\[
u(x,0)=2\sin(3x),
\]

and

\[
u_t(x,0)=0.
\]

Here,

\[
c=2.
\]

The spatial frequency is

\[
n=3.
\]

Therefore the temporal frequency is

\[
cn=6.
\]

Hence,

\begin{answerbox}
\[
\boxed{
u(x,t)
=
2\cos(6t)\sin(3x).
}
\]
\end{answerbox}

\end{workedexamplebox}

%%%%%%%%%%%%%%%%%%%%%%%%%%%%%%%%%%%%%%%%%%%%%%%%%%%%%%%%%%%%
\subsection{Standing Waves}
\label{subsec:standing-waves}
%%%%%%%%%%%%%%%%%%%%%%%%%%%%%%%%%%%%%%%%%%%%%%%%%%%%%%%%%%%%

Separated wave solutions such as

\[
\boxed{
u_n(x,t)
=
\sin
\left(
\frac{n\pi x}{L}
\right)
\cos
\left(
\frac{cn\pi t}{L}
\right)
}
\]

are standing waves.

The spatial nodes remain fixed while the amplitude oscillates in time.

%%%%%%%%%%%%%%%%%%%%%%%%%%%%%%%%%%%%%%%%%%%%%%%%%%%%%%%%%%%%
\subsection{Natural Frequencies}
\label{subsec:wave-natural-frequencies}
%%%%%%%%%%%%%%%%%%%%%%%%%%%%%%%%%%%%%%%%%%%%%%%%%%%%%%%%%%%%

For a string with fixed ends, the angular frequencies are

\[
\boxed{
\omega_n
=
\frac{
cn\pi
}{
L
}.
}
\]

The ordinary frequencies are

\[
\boxed{
f_n
=
\frac{
cn
}{
2L
}.
}
\]

The modes are harmonics of the fundamental frequency.

%%%%%%%%%%%%%%%%%%%%%%%%%%%%%%%%%%%%%%%%%%%%%%%%%%%%%%%%%%%%
\subsection{d'Alembert's Formula}
\label{subsec:dalembert-formula}
%%%%%%%%%%%%%%%%%%%%%%%%%%%%%%%%%%%%%%%%%%%%%%%%%%%%%%%%%%%%

For the wave equation on the whole real line,

\[
u_{tt}
=
c^2u_{xx},
\]

with

\[
u(x,0)=f(x),
\qquad
u_t(x,0)=g(x),
\]

the solution is

\[
\boxed{
u(x,t)
=
\frac12
[
f(x-ct)+f(x+ct)
]
+
\frac1{2c}
\int_{x-ct}^{x+ct}
g(\xi)\,d\xi.
}
\]

%%%%%%%%%%%%%%%%%%%%%%%%%%%%%%%%%%%%%%%%%%%%%%%%%%%%%%%%%%%%
\subsection{Traveling Waves}
\label{subsec:traveling-waves}
%%%%%%%%%%%%%%%%%%%%%%%%%%%%%%%%%%%%%%%%%%%%%%%%%%%%%%%%%%%%

Functions of the form

\[
\boxed{
F(x-ct)
}
\]

move to the right with speed \(c\).

Functions of the form

\[
\boxed{
G(x+ct)
}
\]

move to the left with speed \(c\).

Thus the general one-dimensional homogeneous wave equation naturally
contains two traveling directions.

%%%%%%%%%%%%%%%%%%%%%%%%%%%%%%%%%%%%%%%%%%%%%%%%%%%%%%%%%%%%
\subsection{Finite Propagation Speed}
\label{subsec:finite-propagation-speed}
%%%%%%%%%%%%%%%%%%%%%%%%%%%%%%%%%%%%%%%%%%%%%%%%%%%%%%%%%%%%

From d'Alembert's formula, the value of \(u(x,t)\) depends only on initial
data in the interval

\[
\boxed{
[x-ct,x+ct].
}
\]

Thus disturbances propagate at finite speed \(c\).

This is fundamentally different from the classical heat equation.

%%%%%%%%%%%%%%%%%%%%%%%%%%%%%%%%%%%%%%%%%%%%%%%%%%%%%%%%%%%%
\subsection{Infinite Propagation in Diffusion}
\label{subsec:infinite-propagation-heat}
%%%%%%%%%%%%%%%%%%%%%%%%%%%%%%%%%%%%%%%%%%%%%%%%%%%%%%%%%%%%

For the classical heat equation, localized initial data generally affect
every spatial point immediately for every \(t>0\), although the effect may
be extremely small far away.

Thus the heat equation has mathematically infinite propagation speed.

This reflects the idealized diffusion model.

%%%%%%%%%%%%%%%%%%%%%%%%%%%%%%%%%%%%%%%%%%%%%%%%%%%%%%%%%%%%
\subsection{Wave Energy}
\label{subsec:wave-energy}
%%%%%%%%%%%%%%%%%%%%%%%%%%%%%%%%%%%%%%%%%%%%%%%%%%%%%%%%%%%%

For the fixed-end wave equation,

\[
u_{tt}=c^2u_{xx},
\]

a standard energy is

\[
\boxed{
E(t)
=
\frac12
\int_0^L
\left[
u_t^2
+
c^2u_x^2
\right]
dx.
}
\]

Under appropriate homogeneous boundary conditions, this energy is
conserved.

%%%%%%%%%%%%%%%%%%%%%%%%%%%%%%%%%%%%%%%%%%%%%%%%%%%%%%%%%%%%
\subsection{Energy Conservation for the Wave Equation}
\label{subsec:wave-energy-conservation}
%%%%%%%%%%%%%%%%%%%%%%%%%%%%%%%%%%%%%%%%%%%%%%%%%%%%%%%%%%%%

Differentiate:

\[
E'(t)
=
\int_0^L
\left[
u_tu_{tt}
+
c^2u_xu_{xt}
\right]
dx.
\]

Use

\[
u_{tt}=c^2u_{xx}
\]

and integrate by parts:

\[
E'(t)
=
c^2
\left[
u_tu_x
\right]_0^L.
\]

For fixed endpoints,

\[
u(0,t)=u(L,t)=0,
\]

so

\[
u_t(0,t)=u_t(L,t)=0.
\]

Therefore,

\[
\boxed{
E'(t)=0.
}
\]

%%%%%%%%%%%%%%%%%%%%%%%%%%%%%%%%%%%%%%%%%%%%%%%%%%%%%%%%%%%%
\subsection{Laplace's Equation on a Rectangle}
\label{subsec:laplace-rectangle}
%%%%%%%%%%%%%%%%%%%%%%%%%%%%%%%%%%%%%%%%%%%%%%%%%%%%%%%%%%%%

Consider

\[
\boxed{
u_{xx}+u_{yy}=0
}
\]

on

\[
0<x<L,
\qquad
0<y<H.
\]

Suppose three boundaries are held at zero:

\[
u(0,y)=0,
\]

\[
u(L,y)=0,
\]

\[
u(x,0)=0,
\]

while

\[
u(x,H)=f(x).
\]

%%%%%%%%%%%%%%%%%%%%%%%%%%%%%%%%%%%%%%%%%%%%%%%%%%%%%%%%%%%%
\subsection{Separation for Laplace's Equation}
\label{subsec:laplace-separation}
%%%%%%%%%%%%%%%%%%%%%%%%%%%%%%%%%%%%%%%%%%%%%%%%%%%%%%%%%%%%

Assume

\[
u(x,y)=X(x)Y(y).
\]

Then

\[
X''Y+XY''=0.
\]

Divide by \(XY\):

\[
\frac{X''}{X}
=
-
\frac{Y''}{Y}.
\]

Set both equal to

\[
-\lambda.
\]

Then

\[
\boxed{
X''+\lambda X=0,
}
\]

and

\[
\boxed{
Y''-\lambda Y=0.
}
\]

%%%%%%%%%%%%%%%%%%%%%%%%%%%%%%%%%%%%%%%%%%%%%%%%%%%%%%%%%%%%
\subsection{Rectangle Eigenfunctions}
\label{subsec:rectangle-eigenfunctions}
%%%%%%%%%%%%%%%%%%%%%%%%%%%%%%%%%%%%%%%%%%%%%%%%%%%%%%%%%%%%

The zero conditions at \(x=0,L\) give

\[
\boxed{
X_n(x)
=
\sin
\left(
\frac{n\pi x}{L}
\right).
}
\]

The corresponding \(y\)-equation gives hyperbolic functions.

With \(u(x,0)=0\), the useful solution is

\[
\boxed{
Y_n(y)
=
\sinh
\left(
\frac{n\pi y}{L}
\right).
}
\]

%%%%%%%%%%%%%%%%%%%%%%%%%%%%%%%%%%%%%%%%%%%%%%%%%%%%%%%%%%%%
\subsection{Laplace Rectangle Series}
\label{subsec:laplace-rectangle-series}
%%%%%%%%%%%%%%%%%%%%%%%%%%%%%%%%%%%%%%%%%%%%%%%%%%%%%%%%%%%%

The solution has form

\[
\boxed{
u(x,y)
=
\sum_{n=1}^{\infty}
A_n
\sinh
\left(
\frac{n\pi y}{L}
\right)
\sin
\left(
\frac{n\pi x}{L}
\right).
}
\]

Imposing

\[
u(x,H)=f(x)
\]

gives

\[
A_n
\sinh
\left(
\frac{n\pi H}{L}
\right)
=
b_n,
\]

where \(b_n\) are the sine-series coefficients of \(f\).

Thus,

\[
\boxed{
A_n
=
\frac{
b_n
}{
\sinh
\left(
\frac{n\pi H}{L}
\right)
}.
}
\]

%%%%%%%%%%%%%%%%%%%%%%%%%%%%%%%%%%%%%%%%%%%%%%%%%%%%%%%%%%%%
\subsection{Harmonic Functions}
\label{subsec:harmonic-functions}
%%%%%%%%%%%%%%%%%%%%%%%%%%%%%%%%%%%%%%%%%%%%%%%%%%%%%%%%%%%%

A twice-differentiable function \(u\) is harmonic on a domain if

\[
\boxed{
\Delta u=0.
}
\]

Harmonic functions have exceptionally strong regularity and averaging
properties.

They are central to:

\[
\boxed{
\text{potential theory},
}
\]

\[
\boxed{
\text{complex analysis},
}
\]

and

\[
\boxed{
\text{elliptic PDE theory}.
}
\]

%%%%%%%%%%%%%%%%%%%%%%%%%%%%%%%%%%%%%%%%%%%%%%%%%%%%%%%%%%%%
\subsection{Mean-Value Property Preview}
\label{subsec:harmonic-mean-value-property}
%%%%%%%%%%%%%%%%%%%%%%%%%%%%%%%%%%%%%%%%%%%%%%%%%%%%%%%%%%%%

A harmonic function satisfies a mean-value property.

For an appropriate ball \(B_r(x_0)\) contained in the domain,

\[
\boxed{
u(x_0)
=
\frac1{
|B_r|
}
\int_{B_r(x_0)}
u(x)\,dx.
}
\]

There is also an equivalent spherical-boundary mean-value formulation.

Thus the value at the center is determined by surrounding values.

%%%%%%%%%%%%%%%%%%%%%%%%%%%%%%%%%%%%%%%%%%%%%%%%%%%%%%%%%%%%
\subsection{Maximum Principle}
\label{subsec:maximum-principle}
%%%%%%%%%%%%%%%%%%%%%%%%%%%%%%%%%%%%%%%%%%%%%%%%%%%%%%%%%%%%

A nonconstant harmonic function cannot attain a strict interior maximum
or minimum.

For a bounded domain under suitable hypotheses, its extrema occur on the
boundary.

This is the maximum principle.

%%%%%%%%%%%%%%%%%%%%%%%%%%%%%%%%%%%%%%%%%%%%%%%%%%%%%%%%%%%%
\subsection{Uniqueness from the Maximum Principle}
\label{subsec:laplace-uniqueness-maximum-principle}
%%%%%%%%%%%%%%%%%%%%%%%%%%%%%%%%%%%%%%%%%%%%%%%%%%%%%%%%%%%%

Suppose \(u\) and \(v\) solve the same Dirichlet problem for Laplace's
equation.

Then

\[
w=u-v
\]

satisfies

\[
\Delta w=0
\]

and

\[
w=0
\]

on the boundary.

By the maximum principle,

\[
\boxed{
w=0
}
\]

throughout the domain.

Therefore,

\[
\boxed{
u=v.
}
\]

%%%%%%%%%%%%%%%%%%%%%%%%%%%%%%%%%%%%%%%%%%%%%%%%%%%%%%%%%%%%
\subsection{PDE Existence and Uniqueness}
\label{subsec:pde-existence-uniqueness}
%%%%%%%%%%%%%%%%%%%%%%%%%%%%%%%%%%%%%%%%%%%%%%%%%%%%%%%%%%%%

As with ODEs, two central questions are:

\[
\boxed{
\text{Does a solution exist?}
}
\]

and

\[
\boxed{
\text{Is it unique?}
}
\]

For PDEs, the answer depends strongly on:

\[
\boxed{
\text{PDE type},
}
\]

\[
\boxed{
\text{domain},
}
\]

\[
\boxed{
\text{boundary conditions},
}
\]

and

\[
\boxed{
\text{regularity of the data}.
}
\]

%%%%%%%%%%%%%%%%%%%%%%%%%%%%%%%%%%%%%%%%%%%%%%%%%%%%%%%%%%%%
\subsection{Well-Posed Problems}
\label{subsec:well-posed-pde}
%%%%%%%%%%%%%%%%%%%%%%%%%%%%%%%%%%%%%%%%%%%%%%%%%%%%%%%%%%%%

A problem is called well posed in the sense of Hadamard when:

\[
\boxed{
\text{a solution exists},
}
\]

\[
\boxed{
\text{the solution is unique},
}
\]

and

\[
\boxed{
\text{the solution depends continuously on the data}.
}
\]

Failure of continuous dependence can make a problem unstable in practice.

%%%%%%%%%%%%%%%%%%%%%%%%%%%%%%%%%%%%%%%%%%%%%%%%%%%%%%%%%%%%
\subsection{Fundamental Solution of the Heat Equation}
\label{subsec:heat-fundamental-solution}
%%%%%%%%%%%%%%%%%%%%%%%%%%%%%%%%%%%%%%%%%%%%%%%%%%%%%%%%%%%%

On the real line, the heat equation

\[
u_t=\alpha^2u_{xx}
\]

has fundamental solution

\[
\boxed{
G(x,t)
=
\frac1{
\sqrt{
4\pi\alpha^2t
}
}
\exp
\left(
-
\frac{x^2}{
4\alpha^2t
}
\right),
\qquad
t>0.
}
\]

This is a Gaussian density in \(x\).

%%%%%%%%%%%%%%%%%%%%%%%%%%%%%%%%%%%%%%%%%%%%%%%%%%%%%%%%%%%%
\subsection{Heat Solution by Convolution}
\label{subsec:heat-convolution-solution}
%%%%%%%%%%%%%%%%%%%%%%%%%%%%%%%%%%%%%%%%%%%%%%%%%%%%%%%%%%%%

For initial data

\[
u(x,0)=f(x),
\]

the whole-line solution is

\[
\boxed{
u(x,t)
=
\int_{-\infty}^{\infty}
G(x-\xi,t)
f(\xi)
\,d\xi.
}
\]

Thus,

\[
\boxed{
u(\cdot,t)
=
G_t*f.
}
\]

The Gaussian kernel smooths the initial data.

%%%%%%%%%%%%%%%%%%%%%%%%%%%%%%%%%%%%%%%%%%%%%%%%%%%%%%%%%%%%
\subsection{Probability Connection to the Heat Equation}
\label{subsec:heat-probability-connection}
%%%%%%%%%%%%%%%%%%%%%%%%%%%%%%%%%%%%%%%%%%%%%%%%%%%%%%%%%%%%

The heat kernel has the form of a normal probability density.

This connects diffusion with Brownian motion.

Under an appropriate convention, the density of Brownian motion satisfies
a heat-type equation.

Thus parabolic PDEs have deep connections with stochastic processes.

%%%%%%%%%%%%%%%%%%%%%%%%%%%%%%%%%%%%%%%%%%%%%%%%%%%%%%%%%%%%
\subsection{Green's Functions}
\label{subsec:pde-greens-functions}
%%%%%%%%%%%%%%%%%%%%%%%%%%%%%%%%%%%%%%%%%%%%%%%%%%%%%%%%%%%%

A Green's function represents the response of a linear differential
operator to a point source.

Symbolically, if

\[
L
\]

is a linear operator, a Green's function satisfies

\[
\boxed{
L G(\cdot,\xi)
=
\delta(\cdot-\xi),
}
\]

together with the relevant boundary conditions.

%%%%%%%%%%%%%%%%%%%%%%%%%%%%%%%%%%%%%%%%%%%%%%%%%%%%%%%%%%%%
\subsection{Solution Representation Using a Green's Function}
\label{subsec:greens-function-representation}
%%%%%%%%%%%%%%%%%%%%%%%%%%%%%%%%%%%%%%%%%%%%%%%%%%%%%%%%%%%%

For a linear equation

\[
L u=f,
\]

a Green's function can produce an integral representation such as

\[
\boxed{
u(x)
=
\int
G(x,\xi)
f(\xi)
\,d\xi
}
\]

plus any terms required by boundary data.

This is the PDE analogue of convolution and impulse-response methods.

%%%%%%%%%%%%%%%%%%%%%%%%%%%%%%%%%%%%%%%%%%%%%%%%%%%%%%%%%%%%
\subsection{Fourier Transform Solution of the Heat Equation}
\label{subsec:fourier-heat-equation}
%%%%%%%%%%%%%%%%%%%%%%%%%%%%%%%%%%%%%%%%%%%%%%%%%%%%%%%%%%%%

On the real line,

\[
u_t
=
\alpha^2u_{xx}.
\]

Take the Fourier transform in \(x\):

\[
\boxed{
\widehat{u}_t
=
-\alpha^2\omega^2
\widehat{u}.
}
\]

For each \(\omega\), this is an ODE in time.

Thus,

\[
\boxed{
\widehat{u}(\omega,t)
=
e^{-\alpha^2\omega^2t}
\widehat{f}(\omega).
}
\]

Inverse transformation recovers \(u(x,t)\).

%%%%%%%%%%%%%%%%%%%%%%%%%%%%%%%%%%%%%%%%%%%%%%%%%%%%%%%%%%%%
\subsection{Fourier Transform Solution of the Wave Equation}
\label{subsec:fourier-wave-equation}
%%%%%%%%%%%%%%%%%%%%%%%%%%%%%%%%%%%%%%%%%%%%%%%%%%%%%%%%%%%%

For

\[
u_{tt}
=
c^2u_{xx},
\]

the Fourier transform in \(x\) gives

\[
\boxed{
\widehat{u}_{tt}
+
c^2\omega^2
\widehat{u}
=
0.
}
\]

Thus each spatial frequency behaves like a harmonic oscillator.

This gives

\[
\boxed{
\widehat{u}(\omega,t)
=
A(\omega)\cos(c|\omega|t)
+
B(\omega)\sin(c|\omega|t).
}
\]

%%%%%%%%%%%%%%%%%%%%%%%%%%%%%%%%%%%%%%%%%%%%%%%%%%%%%%%%%%%%
\subsection{Laplace Transform in Time}
\label{subsec:pde-laplace-time-transform}
%%%%%%%%%%%%%%%%%%%%%%%%%%%%%%%%%%%%%%%%%%%%%%%%%%%%%%%%%%%%

For time-dependent PDEs, one can take the Laplace transform in \(t\).

For example,

\[
u_t
=
\alpha^2u_{xx}
\]

becomes

\[
\boxed{
sU(x,s)
-
u(x,0)
=
\alpha^2U_{xx}(x,s).
}
\]

The PDE becomes an ODE in \(x\).

Thus transform methods reduce the number of independent variables.

%%%%%%%%%%%%%%%%%%%%%%%%%%%%%%%%%%%%%%%%%%%%%%%%%%%%%%%%%%%%
\subsection{Conservation Laws}
\label{subsec:conservation-laws-pde}
%%%%%%%%%%%%%%%%%%%%%%%%%%%%%%%%%%%%%%%%%%%%%%%%%%%%%%%%%%%%

A one-dimensional conservation law has form

\[
\boxed{
u_t
+
[f(u)]_x
=
0.
}
\]

Here,

\[
u(x,t)
\]

is a conserved density and

\[
f(u)
\]

is the flux.

This equation expresses local conservation.

%%%%%%%%%%%%%%%%%%%%%%%%%%%%%%%%%%%%%%%%%%%%%%%%%%%%%%%%%%%%
\subsection{Integral Form of a Conservation Law}
\label{subsec:integral-conservation-law}
%%%%%%%%%%%%%%%%%%%%%%%%%%%%%%%%%%%%%%%%%%%%%%%%%%%%%%%%%%%%

Integrate over an interval \([a,b]\):

\[
\frac{d}{dt}
\int_a^b
u(x,t)\,dx
=
-
\int_a^b
[f(u)]_x\,dx.
\]

Thus,

\[
\boxed{
\frac{d}{dt}
\int_a^b
u(x,t)\,dx
=
f(u(a,t))
-
f(u(b,t)).
}
\]

The amount changes only through boundary flux.

%%%%%%%%%%%%%%%%%%%%%%%%%%%%%%%%%%%%%%%%%%%%%%%%%%%%%%%%%%%%
\subsection{Inviscid Burgers Equation}
\label{subsec:inviscid-burgers}
%%%%%%%%%%%%%%%%%%%%%%%%%%%%%%%%%%%%%%%%%%%%%%%%%%%%%%%%%%%%

A fundamental nonlinear conservation law is

\[
\boxed{
u_t
+
u u_x
=
0.
}
\]

Since

\[
u u_x
=
\left(
\frac12u^2
\right)_x,
\]

it can be written

\[
\boxed{
u_t
+
\left(
\frac12u^2
\right)_x
=
0.
}
\]

%%%%%%%%%%%%%%%%%%%%%%%%%%%%%%%%%%%%%%%%%%%%%%%%%%%%%%%%%%%%
\subsection{Characteristics for Burgers Equation}
\label{subsec:burgers-characteristics}
%%%%%%%%%%%%%%%%%%%%%%%%%%%%%%%%%%%%%%%%%%%%%%%%%%%%%%%%%%%%

Along a characteristic,

\[
\boxed{
\frac{dx}{dt}=u.
}
\]

Meanwhile,

\[
\boxed{
\frac{du}{dt}=0.
}
\]

Thus \(u\) is constant along each characteristic, and the characteristic
speed equals the solution value itself.

Different characteristics can therefore move at different speeds.

%%%%%%%%%%%%%%%%%%%%%%%%%%%%%%%%%%%%%%%%%%%%%%%%%%%%%%%%%%%%
\subsection{Shock Formation}
\label{subsec:shock-formation}
%%%%%%%%%%%%%%%%%%%%%%%%%%%%%%%%%%%%%%%%%%%%%%%%%%%%%%%%%%%%

If faster characteristics begin behind slower ones, they can intersect.

At such an intersection, the classical differentiable solution becomes
multivalued if continued naively.

Physically, the solution develops a discontinuity called a shock.

Thus smooth initial data can develop nonsmooth solutions in finite time.

%%%%%%%%%%%%%%%%%%%%%%%%%%%%%%%%%%%%%%%%%%%%%%%%%%%%%%%%%%%%
\subsection{Weak Solutions}
\label{subsec:weak-solutions-preview}
%%%%%%%%%%%%%%%%%%%%%%%%%%%%%%%%%%%%%%%%%%%%%%%%%%%%%%%%%%%%

A weak solution satisfies the PDE after integration against suitable test
functions rather than pointwise differentiation.

This allows discontinuous solutions such as shocks.

Weak formulations are fundamental in modern PDE theory.

%%%%%%%%%%%%%%%%%%%%%%%%%%%%%%%%%%%%%%%%%%%%%%%%%%%%%%%%%%%%
\subsection{Rankine--Hugoniot Condition Preview}
\label{subsec:rankine-hugoniot-preview}
%%%%%%%%%%%%%%%%%%%%%%%%%%%%%%%%%%%%%%%%%%%%%%%%%%%%%%%%%%%%

For

\[
u_t+[f(u)]_x=0,
\]

a shock connecting states

\[
u_L
\]

and

\[
u_R
\]

moves with speed

\[
\boxed{
s
=
\frac{
f(u_L)-f(u_R)
}{
u_L-u_R
}.
}
\]

This is the Rankine--Hugoniot condition.

%%%%%%%%%%%%%%%%%%%%%%%%%%%%%%%%%%%%%%%%%%%%%%%%%%%%%%%%%%%%
\subsection{Entropy Conditions Preview}
\label{subsec:entropy-conditions-preview}
%%%%%%%%%%%%%%%%%%%%%%%%%%%%%%%%%%%%%%%%%%%%%%%%%%%%%%%%%%%%

A conservation law may admit many weak solutions.

Additional admissibility criteria called entropy conditions are used to
select the physically meaningful solution.

These criteria encode the direction in which information and
dissipation should propagate.

%%%%%%%%%%%%%%%%%%%%%%%%%%%%%%%%%%%%%%%%%%%%%%%%%%%%%%%%%%%%
\subsection{Reaction--Diffusion Equations}
\label{subsec:reaction-diffusion}
%%%%%%%%%%%%%%%%%%%%%%%%%%%%%%%%%%%%%%%%%%%%%%%%%%%%%%%%%%%%

A reaction--diffusion equation combines local nonlinear dynamics with
spatial diffusion:

\[
\boxed{
u_t
=
D\Delta u
+
f(u).
}
\]

Here,

\[
D>0
\]

controls diffusion and

\[
f(u)
\]

describes local reaction or growth.

%%%%%%%%%%%%%%%%%%%%%%%%%%%%%%%%%%%%%%%%%%%%%%%%%%%%%%%%%%%%
\subsection{Fisher--KPP Equation Preview}
\label{subsec:fisher-kpp-preview}
%%%%%%%%%%%%%%%%%%%%%%%%%%%%%%%%%%%%%%%%%%%%%%%%%%%%%%%%%%%%

A classical reaction--diffusion equation is

\[
\boxed{
u_t
=
D u_{xx}
+
ru
\left(
1-\frac{u}{K}
\right).
}
\]

It combines logistic growth with spatial diffusion.

The equation supports traveling-wave solutions connecting different
equilibrium states.

%%%%%%%%%%%%%%%%%%%%%%%%%%%%%%%%%%%%%%%%%%%%%%%%%%%%%%%%%%%%
\subsection{Systems of PDEs}
\label{subsec:systems-of-pdes}
%%%%%%%%%%%%%%%%%%%%%%%%%%%%%%%%%%%%%%%%%%%%%%%%%%%%%%%%%%%%

Many physical models involve several interacting fields.

A system may be written

\[
\boxed{
\mathbf{u}_t
=
D\Delta\mathbf{u}
+
\mathbf{f}(\mathbf{u}).
}
\]

Examples include:

\[
\boxed{
\text{reaction--diffusion systems},
}
\]

\[
\boxed{
\text{fluid equations},
}
\]

and

\[
\boxed{
\text{electromagnetic systems}.
}
\]

%%%%%%%%%%%%%%%%%%%%%%%%%%%%%%%%%%%%%%%%%%%%%%%%%%%%%%%%%%%%
\subsection{Elliptic, Parabolic, and Hyperbolic Behavior}
\label{subsec:pde-type-comparison}
%%%%%%%%%%%%%%%%%%%%%%%%%%%%%%%%%%%%%%%%%%%%%%%%%%%%%%%%%%%%

A useful qualitative comparison is:

\[
\boxed{
\text{elliptic}
\rightarrow
\text{equilibrium and spatial smoothing},
}
\]

\[
\boxed{
\text{parabolic}
\rightarrow
\text{diffusion and temporal smoothing},
}
\]

\[
\boxed{
\text{hyperbolic}
\rightarrow
\text{wave propagation and characteristics}.
}
\]

These are guiding interpretations rather than universal descriptions.

%%%%%%%%%%%%%%%%%%%%%%%%%%%%%%%%%%%%%%%%%%%%%%%%%%%%%%%%%%%%
\subsection{Separation of Variables Workflow}
\label{subsec:pde-separation-workflow}
%%%%%%%%%%%%%%%%%%%%%%%%%%%%%%%%%%%%%%%%%%%%%%%%%%%%%%%%%%%%

For a separable linear boundary-value problem:

\begin{enumerate}
    \item assume a product solution such as
          \(u(x,t)=X(x)T(t)\);
    \item substitute into the PDE;
    \item divide by the product;
    \item set the separated expressions equal to a constant;
    \item solve the spatial eigenvalue problem;
    \item solve the temporal equation for each eigenvalue;
    \item superpose all separated modes;
    \item expand the initial or boundary data in the corresponding
          eigenfunction series;
    \item determine the coefficients by orthogonality.
\end{enumerate}

%%%%%%%%%%%%%%%%%%%%%%%%%%%%%%%%%%%%%%%%%%%%%%%%%%%%%%%%%%%%
\subsection{PDE Modeling Workflow}
\label{subsec:pde-modeling-workflow}
%%%%%%%%%%%%%%%%%%%%%%%%%%%%%%%%%%%%%%%%%%%%%%%%%%%%%%%%%%%%

A practical PDE modeling workflow is:

\begin{enumerate}
    \item identify the dependent field;
    \item identify spatial and temporal variables;
    \item derive or state the governing balance law;
    \item determine the spatial domain;
    \item specify initial conditions;
    \item specify boundary conditions;
    \item classify the PDE;
    \item determine whether symmetry reduces the problem;
    \item identify applicable transform or separation methods;
    \item determine whether a closed-form solution is realistic;
    \item analyze qualitative properties;
    \item use numerical methods when necessary;
    \item check dimensional consistency;
    \item interpret the solution physically.
\end{enumerate}

%%%%%%%%%%%%%%%%%%%%%%%%%%%%%%%%%%%%%%%%%%%%%%%%%%%%%%%%%%%%
\subsection{Numerical PDE Methods Preview}
\label{subsec:numerical-pde-preview}
%%%%%%%%%%%%%%%%%%%%%%%%%%%%%%%%%%%%%%%%%%%%%%%%%%%%%%%%%%%%

Most realistic PDEs do not admit elementary closed-form solutions.

Common numerical approaches include:

\[
\boxed{
\text{finite differences},
}
\]

\[
\boxed{
\text{finite elements},
}
\]

\[
\boxed{
\text{finite volumes},
}
\]

and

\[
\boxed{
\text{spectral methods}.
}
\]

These are developed more deeply in Chapter 11.

%%%%%%%%%%%%%%%%%%%%%%%%%%%%%%%%%%%%%%%%%%%%%%%%%%%%%%%%%%%%
\subsection{Finite-Difference Approximation}
\label{subsec:finite-difference-pde}
%%%%%%%%%%%%%%%%%%%%%%%%%%%%%%%%%%%%%%%%%%%%%%%%%%%%%%%%%%%%

A centered approximation to the second derivative is

\[
\boxed{
u_{xx}(x_i,t)
\approx
\frac{
u_{i+1}-2u_i+u_{i-1}
}{
(\Delta x)^2
}.
}
\]

For the time derivative,

\[
\boxed{
u_t
\approx
\frac{
u_i^{n+1}-u_i^n
}{
\Delta t
}.
}
\]

Combining these gives a discrete approximation to the heat equation.

%%%%%%%%%%%%%%%%%%%%%%%%%%%%%%%%%%%%%%%%%%%%%%%%%%%%%%%%%%%%
\subsection{Explicit Heat Scheme Preview}
\label{subsec:explicit-heat-scheme}
%%%%%%%%%%%%%%%%%%%%%%%%%%%%%%%%%%%%%%%%%%%%%%%%%%%%%%%%%%%%

For

\[
u_t=\alpha^2u_{xx},
\]

an explicit scheme is

\[
\boxed{
u_i^{n+1}
=
u_i^n
+
r
\left(
u_{i+1}^n
-
2u_i^n
+
u_{i-1}^n
\right),
}
\]

where

\[
\boxed{
r
=
\frac{
\alpha^2\Delta t
}{
(\Delta x)^2
}.
}
\]

%%%%%%%%%%%%%%%%%%%%%%%%%%%%%%%%%%%%%%%%%%%%%%%%%%%%%%%%%%%%
\subsection{Stability Restriction Preview}
\label{subsec:heat-cfl-preview}
%%%%%%%%%%%%%%%%%%%%%%%%%%%%%%%%%%%%%%%%%%%%%%%%%%%%%%%%%%%%

For the standard one-dimensional explicit heat scheme, stability requires

\[
\boxed{
r
\leq
\frac12.
}
\]

Thus

\[
\boxed{
\Delta t
\leq
\frac{
(\Delta x)^2
}{
2\alpha^2
}.
}
\]

A numerically convenient discretization can therefore still be unstable
if the time step is too large.

%%%%%%%%%%%%%%%%%%%%%%%%%%%%%%%%%%%%%%%%%%%%%%%%%%%%%%%%%%%%
\subsection{Finite-Element Method Preview}
\label{subsec:finite-element-preview}
%%%%%%%%%%%%%%%%%%%%%%%%%%%%%%%%%%%%%%%%%%%%%%%%%%%%%%%%%%%%

Finite-element methods begin from a weak formulation of the PDE.

The domain is divided into elements, and the approximate solution is
constructed from basis functions.

This approach is especially powerful for:

\[
\boxed{
\text{irregular geometries},
}
\]

\[
\boxed{
\text{variable coefficients},
}
\]

and

\[
\boxed{
\text{complex boundary conditions}.
}
\]

%%%%%%%%%%%%%%%%%%%%%%%%%%%%%%%%%%%%%%%%%%%%%%%%%%%%%%%%%%%%
\subsection{Spectral Methods Preview}
\label{subsec:spectral-methods-pde-preview}
%%%%%%%%%%%%%%%%%%%%%%%%%%%%%%%%%%%%%%%%%%%%%%%%%%%%%%%%%%%%

Spectral methods approximate the solution using global basis functions
such as:

\[
\boxed{
\text{Fourier modes},
}
\]

or

\[
\boxed{
\text{orthogonal polynomials}.
}
\]

For sufficiently smooth solutions, spectral methods can achieve very high
accuracy with relatively few coefficients.

%%%%%%%%%%%%%%%%%%%%%%%%%%%%%%%%%%%%%%%%%%%%%%%%%%%%%%%%%%%%
\subsection{Common Errors}
\label{subsec:pde-common-errors}
%%%%%%%%%%%%%%%%%%%%%%%%%%%%%%%%%%%%%%%%%%%%%%%%%%%%%%%%%%%%

\subsubsection{Ignoring Boundary Conditions}

A PDE alone usually does not determine a unique solution.

Boundary and initial conditions are part of the mathematical problem.

%%%%%%%%%%%%%%%%%%%%%%%%%%%%%%%%%%%%%%%%%%%%%%%%%%%%%%%%%%%%
\subsubsection{Using Too Few Initial Conditions for a Wave Equation}

Because the wave equation is second order in time, one generally needs
both

\[
u(x,0)
\]

and

\[
u_t(x,0).
\]

%%%%%%%%%%%%%%%%%%%%%%%%%%%%%%%%%%%%%%%%%%%%%%%%%%%%%%%%%%%%
\subsubsection{Using Heat-Equation Logic for Wave Problems}

Diffusion equations smooth and dissipate.

Wave equations transport and oscillate.

Their qualitative behavior is fundamentally different.

%%%%%%%%%%%%%%%%%%%%%%%%%%%%%%%%%%%%%%%%%%%%%%%%%%%%%%%%%%%%
\subsubsection{Confusing Dirichlet and Neumann Conditions}

Dirichlet conditions specify

\[
\boxed{
u.
}
\]

Neumann conditions specify

\[
\boxed{
\frac{\partial u}{\partial n}.
}
\]

They model different physical constraints.

%%%%%%%%%%%%%%%%%%%%%%%%%%%%%%%%%%%%%%%%%%%%%%%%%%%%%%%%%%%%
\subsubsection{Forgetting the Factor \(2B\) in PDE Classification}

If the PDE is written

\[
A u_{xx}
+
2B u_{xy}
+
C u_{yy},
\]

the discriminant is

\[
\boxed{
B^2-AC.
}
\]

If the mixed coefficient is written instead as \(Bu_{xy}\), the
classification formula must be adjusted accordingly.

%%%%%%%%%%%%%%%%%%%%%%%%%%%%%%%%%%%%%%%%%%%%%%%%%%%%%%%%%%%%
\subsubsection{Choosing the Wrong Separation Sign}

For fixed-end heat or wave problems, choosing the separation constant
incorrectly can produce only trivial boundary solutions.

The boundary-value problem determines the allowable sign and eigenvalues.

%%%%%%%%%%%%%%%%%%%%%%%%%%%%%%%%%%%%%%%%%%%%%%%%%%%%%%%%%%%%
\subsubsection{Forgetting Eigenfunction Orthogonality}

Fourier coefficients arise from orthogonality.

The normalization factor must match the interval and basis used.

%%%%%%%%%%%%%%%%%%%%%%%%%%%%%%%%%%%%%%%%%%%%%%%%%%%%%%%%%%%%
\subsubsection{Assuming Every PDE Is Separable}

Separation of variables works only for sufficiently structured equations,
domains, and boundary conditions.

Most PDEs are not solvable by elementary separation.

%%%%%%%%%%%%%%%%%%%%%%%%%%%%%%%%%%%%%%%%%%%%%%%%%%%%%%%%%%%%
\subsubsection{Confusing Steady State with Initial State}

Steady state means

\[
\boxed{
u_t=0.
}
\]

It describes time-independent equilibrium, not necessarily the condition
at \(t=0\).

%%%%%%%%%%%%%%%%%%%%%%%%%%%%%%%%%%%%%%%%%%%%%%%%%%%%%%%%%%%%
\subsubsection{Forgetting High-Frequency Heat Modes Decay Faster}

For the heat equation,

\[
e^{-Cn^2t}
\]

shows that large \(n\) modes disappear more rapidly.

This is why diffusion smooths solutions.

%%%%%%%%%%%%%%%%%%%%%%%%%%%%%%%%%%%%%%%%%%%%%%%%%%%%%%%%%%%%
\subsubsection{Assuming Heat Disturbances Travel at Finite Speed}

The classical heat equation produces immediate nonzero influence
throughout space for positive time.

The wave equation instead has finite propagation speed.

%%%%%%%%%%%%%%%%%%%%%%%%%%%%%%%%%%%%%%%%%%%%%%%%%%%%%%%%%%%%
\subsubsection{Treating Shock Solutions as Classical Solutions}

A shock is discontinuous and does not satisfy the conservation law
pointwise in the classical derivative sense.

A weak formulation is required.

%%%%%%%%%%%%%%%%%%%%%%%%%%%%%%%%%%%%%%%%%%%%%%%%%%%%%%%%%%%%
\subsubsection{Assuming Every Weak Solution Is Physically Relevant}

Nonlinear conservation laws may admit multiple weak solutions.

Entropy conditions may be required to select the admissible one.

%%%%%%%%%%%%%%%%%%%%%%%%%%%%%%%%%%%%%%%%%%%%%%%%%%%%%%%%%%%%
\subsubsection{Ignoring Numerical Stability}

A consistent finite-difference method can still fail numerically if the
stability condition is violated.

%%%%%%%%%%%%%%%%%%%%%%%%%%%%%%%%%%%%%%%%%%%%%%%%%%%%%%%%%%%%
\subsection{Applications}
\label{subsec:pde-applications}
%%%%%%%%%%%%%%%%%%%%%%%%%%%%%%%%%%%%%%%%%%%%%%%%%%%%%%%%%%%%

\begin{applicationbox}

\textbf{Heat Transfer.}
The heat equation models conduction and diffusion of thermal energy.

\medskip

\textbf{Wave Motion.}
The wave equation describes vibrating strings, acoustics, elastic media,
and electromagnetic propagation.

\medskip

\textbf{Electrostatics.}
Laplace and Poisson equations describe electric potentials generated by
charge distributions.

\medskip

\textbf{Fluid Mechanics.}
Conservation laws and systems of nonlinear PDEs govern fluid motion.

\medskip

\textbf{Environmental Science.}
Advection--diffusion and reaction--diffusion equations model transport of
pollutants, nutrients, oxygen, and ecological populations.

\medskip

\textbf{Biology.}
Reaction--diffusion systems model spatial population dynamics, pattern
formation, and traveling biological waves.

\medskip

\textbf{Quantum Mechanics.}
The Schrödinger equation is a fundamental PDE describing quantum-state
evolution.

\medskip

\textbf{Probability.}
Diffusion processes are connected with heat-type equations and stochastic
differential equations.

\medskip

\textbf{Geometry.}
Elliptic and nonlinear PDEs appear throughout differential geometry and
geometric analysis.

\end{applicationbox}

%%%%%%%%%%%%%%%%%%%%%%%%%%%%%%%%%%%%%%%%%%%%%%%%%%%%%%%%%%%%
\subsection{Exercises}
\label{subsec:pde-exercises}
%%%%%%%%%%%%%%%%%%%%%%%%%%%%%%%%%%%%%%%%%%%%%%%%%%%%%%%%%%%%

\begin{exercise}[1]
Define a partial differential equation.
\end{exercise}

\begin{exercise}[1]
Determine the order of

\[
u_{xxy}
+
u_x
=
0.
\]
\end{exercise}

\begin{exercise}[1]
Classify

\[
u_{xx}
+
u_{yy}
=
0
\]

as elliptic, parabolic, or hyperbolic.
\end{exercise}

\begin{exercise}[1]
Classify

\[
u_t
=
u_{xx}.
\]
\end{exercise}

\begin{exercise}[1]
Classify

\[
u_{tt}
=
4u_{xx}.
\]
\end{exercise}

\begin{exercise}[1]
Define a Dirichlet boundary condition.
\end{exercise}

\begin{exercise}[1]
Define a Neumann boundary condition.
\end{exercise}

\begin{exercise}[1]
Define a Robin boundary condition.
\end{exercise}

\begin{exercise}[1]
Define a harmonic function.
\end{exercise}

\begin{exercise}[1]
State Laplace's equation.
\end{exercise}

\begin{exercise}[1]
State the heat equation.
\end{exercise}

\begin{exercise}[1]
State the wave equation.
\end{exercise}

\begin{exercise}[2]
Verify that

\[
u(x,t)
=
e^{-t}\sin x
\]

solves

\[
u_t=u_{xx}.
\]
\end{exercise}

\begin{exercise}[2]
Verify that

\[
u(x,t)
=
\cos(2t)\sin x
\]

does or does not solve

\[
u_{tt}=4u_{xx}.
\]
\end{exercise}

\begin{exercise}[2]
Solve

\[
u_t+2u_x=0,
\qquad
u(x,0)=\cos x.
\]
\end{exercise}

\begin{exercise}[2]
Find the steady-state solution satisfying

\[
u_{xx}=0,
\]

\[
u(0)=10,
\qquad
u(5)=30.
\]
\end{exercise}

\begin{exercise}[2]
Find the eigenvalues of

\[
X''+\lambda X=0,
\qquad
X(0)=X(L)=0.
\]
\end{exercise}

\begin{exercise}[2]
Find the corresponding eigenfunctions.
\end{exercise}

\begin{exercise}[2]
Solve

\[
u_t=u_{xx},
\qquad
0<x<\pi,
\]

with zero boundary values and initial condition

\[
u(x,0)=5\sin x.
\]
\end{exercise}

\begin{exercise}[2]
Solve the same heat equation with initial condition

\[
u(x,0)=2\sin(3x).
\]
\end{exercise}

\begin{exercise}[2]
Solve

\[
u_{tt}=u_{xx},
\]

with zero endpoint values,

\[
u(x,0)=\sin(2x),
\qquad
u_t(x,0)=0.
\]
\end{exercise}

\begin{exercise}[3]
Derive the solution

\[
u(x,t)=f(x-ct)
\]

for the transport equation

\[
u_t+cu_x=0.
\]
\end{exercise}

\begin{exercise}[3]
Derive the characteristic equations for a general first-order
quasilinear PDE.
\end{exercise}

\begin{exercise}[3]
Use separation of variables to derive the heat-equation eigenvalue
problem.
\end{exercise}

\begin{exercise}[3]
Derive the Fourier coefficients for the fixed-end heat equation.
\end{exercise}

\begin{exercise}[3]
Explain why all nonzero heat modes decay exponentially.
\end{exercise}

\begin{exercise}[3]
Explain why higher-frequency heat modes decay faster.
\end{exercise}

\begin{exercise}[3]
Derive the fixed-end wave equation series solution.
\end{exercise}

\begin{exercise}[3]
Derive the coefficients \(A_n\) and \(B_n\) from the initial displacement
and velocity.
\end{exercise}

\begin{exercise}[3]
Derive the conserved wave energy.
\end{exercise}

\begin{exercise}[3]
Verify d'Alembert's formula directly for sufficiently smooth data.
\end{exercise}

\begin{exercise}[3]
Explain finite propagation speed using d'Alembert's formula.
\end{exercise}

\begin{exercise}[3]
Separate Laplace's equation on a rectangle and derive the spatial
eigenfunctions.
\end{exercise}

\begin{exercise}[3]
Explain how the maximum principle implies uniqueness for the Dirichlet
problem.
\end{exercise}

\begin{exercise}[3]
Show that

\[
u(x,y)=x^2-y^2
\]

is harmonic.
\end{exercise}

\begin{exercise}[3]
Determine whether

\[
u(x,y)=x^2+y^2
\]

is harmonic.
\end{exercise}

\begin{exercise}[3]
Derive the Fourier-transform form

\[
\widehat{u}_t
=
-\alpha^2\omega^2\widehat{u}
\]

for the heat equation.
\end{exercise}

\begin{exercise}[3]
Use the Fourier transform to derive the Gaussian heat kernel.
\end{exercise}

\begin{exercise}[3]
Derive the integral conservation law from

\[
u_t+[f(u)]_x=0.
\]
\end{exercise}

\begin{exercise}[3]
Derive the characteristics of Burgers equation.
\end{exercise}

\begin{exercise}[3]
Explain geometrically why intersecting characteristics lead to shock
formation.
\end{exercise}

\begin{exercise}[3]
Derive the Rankine--Hugoniot shock speed.
\end{exercise}

\begin{exercise}[3]
Explain why weak solutions are needed for conservation laws.
\end{exercise}

\begin{exercise}[3]
Derive the standard explicit finite-difference scheme for the heat
equation.
\end{exercise}

\begin{exercise}[3]
Explain the meaning of the ratio

\[
r
=
\frac{
\alpha^2\Delta t
}{
(\Delta x)^2
}.
\]
\end{exercise}

\begin{exercise}[4]
Study Sturm--Liouville theory underlying separation of variables.
\end{exercise}

\begin{exercise}[4]
Prove orthogonality of eigenfunctions for a regular Sturm--Liouville
problem.
\end{exercise}

\begin{exercise}[4]
Study convergence of Fourier-series PDE solutions.
\end{exercise}

\begin{exercise}[4]
Prove the maximum principle for harmonic functions.
\end{exercise}

\begin{exercise}[4]
Study the mean-value property of harmonic functions.
\end{exercise}

\begin{exercise}[4]
Derive Green's identities.
\end{exercise}

\begin{exercise}[4]
Construct Green's functions for simple Poisson problems.
\end{exercise}

\begin{exercise}[4]
Study weak formulations of elliptic boundary-value problems.
\end{exercise}

\begin{exercise}[4]
Prove uniqueness for Poisson's equation using an energy method.
\end{exercise}

\begin{exercise}[4]
Study Sobolev spaces as function spaces for weak PDE solutions.
\end{exercise}

\begin{exercise}[4]
Derive the heat semigroup and its fundamental solution.
\end{exercise}

\begin{exercise}[4]
Study the smoothing properties of the heat equation.
\end{exercise}

\begin{exercise}[4]
Study finite-speed propagation for multidimensional wave equations.
\end{exercise}

\begin{exercise}[4]
Derive Kirchhoff-type formulas for higher-dimensional wave equations.
\end{exercise}

\begin{exercise}[4]
Study entropy solutions for scalar conservation laws.
\end{exercise}

\begin{exercise}[4]
Study rarefaction waves for Burgers equation.
\end{exercise}

\begin{exercise}[4]
Analyze traveling-wave solutions of reaction--diffusion equations.
\end{exercise}

\begin{exercise}[4]
Study the Fisher--KPP minimal wave speed.
\end{exercise}

\begin{exercise}[4]
Derive stability conditions for explicit finite-difference heat schemes.
\end{exercise}

\begin{exercise}[4]
Study finite-element weak formulations.
\end{exercise}

\begin{exercise}[4]
Investigate spectral methods for periodic PDEs.
\end{exercise}

%%%%%%%%%%%%%%%%%%%%%%%%%%%%%%%%%%%%%%%%%%%%%%%%%%%%%%%%%%%%
\subsection{Conceptual Summary}
\label{subsec:pde-summary}
%%%%%%%%%%%%%%%%%%%%%%%%%%%%%%%%%%%%%%%%%%%%%%%%%%%%%%%%%%%%

A partial differential equation involves a function of several variables
and its partial derivatives.

A general second-order linear PDE contains a principal part such as

\[
\boxed{
A u_{xx}
+
2B u_{xy}
+
C u_{yy}.
}
\]

Its classification is determined by

\[
\boxed{
B^2-AC.
}
\]

Laplace's equation is

\[
\boxed{
\Delta u=0.
}
\]

Poisson's equation is

\[
\boxed{
\Delta u=f.
}
\]

The heat equation is

\[
\boxed{
u_t=\alpha^2u_{xx}.
}
\]

The wave equation is

\[
\boxed{
u_{tt}=c^2u_{xx}.
}
\]

The transport equation is

\[
\boxed{
u_t+cu_x=0.
}
\]

Separation of variables writes

\[
\boxed{
u(x,t)=X(x)T(t)
}
\]

and converts the PDE into ODE eigenvalue problems.

The fixed-end heat solution is

\[
\boxed{
u(x,t)
=
\sum_{n=1}^{\infty}
b_n
e^{-\alpha^2(n\pi/L)^2t}
\sin
\left(
\frac{n\pi x}{L}
\right).
}
\]

The fixed-end wave solution is

\[
\boxed{
u(x,t)
=
\sum_{n=1}^{\infty}
\left[
A_n\cos
\left(
\frac{cn\pi t}{L}
\right)
+
B_n\sin
\left(
\frac{cn\pi t}{L}
\right)
\right]
\sin
\left(
\frac{n\pi x}{L}
\right).
}
\]

Partial differential equations therefore combine

\[
\boxed{
\text{multivariable calculus}
+
\text{ODEs}
+
\text{Fourier analysis}
+
\text{linear algebra}
+
\text{physical modeling}.
}
\]

%%%%%%%%%%%%%%%%%%%%%%%%%%%%%%%%%%%%%%%%%%%%%%%%%%%%%%%%%%%%
\subsection{Section Connections}
\label{subsec:pde-connections}
%%%%%%%%%%%%%%%%%%%%%%%%%%%%%%%%%%%%%%%%%%%%%%%%%%%%%%%%%%%%

\chapterconnection{
Partial Differential Equations extend the differential-equation framework
from functions of one independent variable to spatially distributed
fields. Separation of variables connects PDEs with eigenvalue problems,
Fourier series, and ordinary differential equations, while transform
methods from Section~\ref{sec:transform-methods} convert PDEs into lower-
dimensional equations. Elliptic equations describe equilibrium,
parabolic equations describe diffusion, and hyperbolic equations describe
wave propagation. Nonlinear conservation laws introduce shocks and weak
solutions. The next section, Difference Equations, develops discrete-time
analogues of differential equations and connects recurrence relations
with discrete dynamical systems and the \(z\)-transform.
}

%%%%%%%%%%%%%%%%%%%%%%%%%%%%%%%%%%%%%%%%%%%%%%%%%%%%%%%%%%%%
% SECTION SOURCE TRACKING
%%%%%%%%%%%%%%%%%%%%%%%%%%%%%%%%%%%%%%%%%%%%%%%%%%%%%%%%%%%%

%------------------------------------------------------------
% 9.4 Partial Differential Equations
%------------------------------------------------------------
%
% Source basis:
%   - Standard introductory PDE theory.
%   - Standard separation-of-variables and Fourier methods.
%   - Standard heat, wave, Laplace, and transport equations.
%
% Used for:
%   - PDE definitions and classification;
%   - order and linearity;
%   - semilinear, quasilinear, and nonlinear equations;
%   - initial and boundary conditions;
%   - Dirichlet, Neumann, Robin, and mixed conditions;
%   - elliptic, parabolic, and hyperbolic classification;
%   - Laplacian;
%   - Laplace and Poisson equations;
%   - heat equation;
%   - wave equation;
%   - transport equation;
%   - characteristics;
%   - separation of variables;
%   - eigenvalue problems;
%   - Fourier-series solutions;
%   - heat equation on finite intervals;
%   - steady-state heat equations;
%   - wave equation on finite intervals;
%   - standing waves and natural frequencies;
%   - d'Alembert formula;
%   - finite propagation speed;
%   - wave energy;
%   - Laplace equation on rectangles;
%   - harmonic functions;
%   - maximum principle;
%   - PDE uniqueness;
%   - well-posedness;
%   - heat kernel;
%   - Green's functions;
%   - Fourier-transform PDE methods;
%   - Laplace-transform PDE methods;
%   - conservation laws;
%   - Burgers equation;
%   - characteristics and shock formation;
%   - weak solutions;
%   - Rankine--Hugoniot condition;
%   - entropy conditions preview;
%   - reaction--diffusion equations;
%   - Fisher--KPP preview;
%   - systems of PDEs;
%   - finite differences;
%   - finite elements preview;
%   - spectral methods preview;
%   - applications and exercises.
%
% Notes:
%   - Difference Equations is developed next in Section 9.5.
%   - Integral Equations, Delay Differential Equations, and
%     Differential-Algebraic Equations follow later in Chapter 9.
%   - Numerical PDE methods are developed more deeply in Chapter 11.
%
%%%%%%%%%%%%%%%%%%%%%%%%%%%%%%%%%%%%%%%%%%%%%%%%%%%%%%%%%%%%